% 2026-10-13 (user): the bucket story check of the feature study (results/feature_mcs_2026-10-08/SUMMARY.md, section
% "The bucket story check"; experiments/feature_mcs_buckets.py).  Every number is a \bk... macro written by
% writeup/make_buckets_tex.py from buckets_{pairwise,stepm,spa,verdicts,gates}.csv, bucket_losses.parquet and
% mcs_sets.csv (every qualitative claim below is asserted there).  No LSTM is among the study's estimators or pools.
% Input at the end of sections/appendix_campaign_16h.tex, where the parked lean-set subsection was (the feature study
% replaces it; MASTER_TABLE_REVIEW_2026-10-01.md items 84 and 91-92).
\ifdefined\bkNSeries\else% AUTO-GENERATED by writeup/make_buckets_tex.py -- do not edit.
% The bucket story check (experiments/feature_mcs_buckets.py), read from
% results/feature_mcs_2026-10-08/buckets_{pairwise,stepm,spa,verdicts,gates}.csv, bucket_losses.parquet
% and mcs_sets.csv.  Point estimates; DM and HAC t significant at |.| > 1.96; StepM at FWER 5 %.
\newcommand{\bkNBuckets}{8}  % buckets
\newcommand{\bkNSeries}{47}  % series of all_features the buckets partition
\newcommand{\bkSizeMom}{7}  % series in bucket moments
\newcommand{\bkSizeLiq}{14}  % series in bucket liquidity
\newcommand{\bkSizeMew}{6}  % series in bucket market_ew
\newcommand{\bkSizeMvw}{6}  % series in bucket market_vw
\newcommand{\bkSizeSent}{3}  % series in bucket sentiment
\newcommand{\bkSizeIv}{3}  % series in bucket implied_vol
\newcommand{\bkSizeVd}{4}  % series in bucket vol_demand
\newcommand{\bkSizeFomc}{4}  % series in bucket fomc
\newcommand{\bkNEst}{6}  % estimators
\newcommand{\bkNDesigns}{18}  % designs per estimator: 8 add-bucket, 8 leave-bucket-out, 2 anchors
\newcommand{\bkDays}{866}  % trade days
\newcommand{\bkBoot}{10{,}000}  % bootstrap replicates (circular block)
\newcommand{\bkBlockA}{5}
\newcommand{\bkBlockB}{10}
\newcommand{\bkBlockC}{21}
\newcommand{\bkFwer}{5}  % StepM family-wise error rate of the verdict, %
\newcommand{\bkSigPct}{5}  % two-sided level of the pairwise DM and HAC t (|.| > 1.96)
\newcommand{\bkNPools}{4}  % pools
\newcommand{\bkNFam}{10}  % families: six estimators and four pools
\newcommand{\bkCells}{30}  % family x loss cells of the verdict
\newcommand{\bkHolds}{none}  % cells where the story holds
\newcommand{\bkAllBeats}{none}  % cells where all features beats HAR + cal. (StepM, every block)
\newcommand{\bkBucketCells}{four}  % cells where one bucket beats HAR + cal. (all market_vw, linear)
\newcommand{\bkBucketCellsWho}{the ridge on winsorized QLIKE, the lasso on QLIKE, the lasso on winsorized QLIKE, the linear pool on winsorized QLIKE}  % those cells
\newcommand{\bkRemoveCells}{three}  % cells where StepM (test C, every block) says removing a bucket hurts
\newcommand{\bkAllDMMin}{\ensuremath{-}1.43}  % most negative DM QLIKE, all features vs HAR + cal. (lgbm)
\newcommand{\bkAllDMMinWho}{LightGBM}
\newcommand{\bkAllTMax}{1.47}  % largest HAC t, all features vs HAR + cal. (lasso)
\newcommand{\bkAllTMaxWho}{lasso}
\newcommand{\bkAllShHigher}{10}  % families where all features trades higher
\newcommand{\bkLassoVwDQ}{\ensuremath{-}3.4}  % lasso + market_vw vs HAR + cal., QLIKE change %
\newcommand{\bkLassoVwDM}{\ensuremath{-}3.75}
\newcommand{\bkLassoVwSpaMax}{0.002}  % lasso, SPA consistent p (test A, QLIKE), largest over blocks
\newcommand{\bkPoolLinVwDM}{\ensuremath{-}2.35}  % linear pool + market_vw vs HAR + cal., DM QLIKE
\newcommand{\bkPoolLinVwDMw}{\ensuremath{-}2.96}  % ... DM winsorized QLIKE
\newcommand{\bkPoolLinAllDM}{\ensuremath{-}0.23}  % linear pool, all features vs HAR + cal., DM QLIKE
\newcommand{\bkPoolLinAllDMw}{\ensuremath{-}0.66}  % ... DM winsorized QLIKE
\newcommand{\bkPoolLinVwOutDM}{\ensuremath{+}2.22}  % linear pool, all features - market_vw vs all features, DM QLIKE
\newcommand{\bkPoolLinVwOutDMw}{\ensuremath{+}3.00}  % ... winsorized
\newcommand{\bkLiqLinDQMin}{8.2}  % linear + liquidity vs HAR + cal., QLIKE change %, smallest
\newcommand{\bkLiqLinDQMax}{15.0}  % ... largest
\newcommand{\bkLiqLinDMMin}{\ensuremath{+}1.96}  % ... DM smallest
\newcommand{\bkLiqLinDMMax}{\ensuremath{+}2.58}  % ... DM largest
\newcommand{\bkXgbBaseSh}{0.95}  % xgb: HAR + cal. Sharpe mid
\newcommand{\bkXgbAllSh}{1.12}  % xgb: all features Sharpe mid
\newcommand{\bkXgbAllT}{0.30}  % xgb: all features vs HAR + cal., HAC t
\newcommand{\bkXgbMomSh}{1.66}  % xgb + moments: Sharpe mid
\newcommand{\bkXgbMomT}{1.86}  % xgb + moments vs HAR + cal.: HAC t
\newcommand{\bkXgbLiqSh}{1.67}  % xgb + liquidity: Sharpe mid
\newcommand{\bkXgbLiqT}{2.05}  % xgb + liquidity vs HAR + cal.: HAC t
\newcommand{\bkXgbMewSh}{1.65}  % xgb + market_ew: Sharpe mid
\newcommand{\bkXgbMewT}{2.39}  % xgb + market_ew vs HAR + cal.: HAC t
\newcommand{\bkRfBaseSh}{0.66}  % rf: HAR + cal. Sharpe mid
\newcommand{\bkRfAllSh}{0.85}  % rf: all features Sharpe mid
\newcommand{\bkRfAllT}{0.34}  % rf: all features vs HAR + cal., HAC t
\newcommand{\bkRfMomSh}{1.47}  % rf + moments: Sharpe mid
\newcommand{\bkRfMomT}{2.35}  % rf + moments vs HAR + cal.: HAC t
\newcommand{\bkRfLiqSh}{1.05}  % rf + liquidity: Sharpe mid
\newcommand{\bkRfLiqT}{1.07}  % rf + liquidity vs HAR + cal.: HAC t
\newcommand{\bkRfMewSh}{1.47}  % rf + market_ew: Sharpe mid
\newcommand{\bkRfMewT}{2.38}  % rf + market_ew vs HAR + cal.: HAC t
\newcommand{\bkLgbmBaseSh}{0.79}  % lgbm: HAR + cal. Sharpe mid
\newcommand{\bkLgbmAllSh}{1.33}  % lgbm: all features Sharpe mid
\newcommand{\bkLgbmAllT}{1.11}  % lgbm: all features vs HAR + cal., HAC t
\newcommand{\bkLgbmMomSh}{1.61}  % lgbm + moments: Sharpe mid
\newcommand{\bkLgbmMomT}{1.48}  % lgbm + moments vs HAR + cal.: HAC t
\newcommand{\bkLgbmLiqSh}{1.16}  % lgbm + liquidity: Sharpe mid
\newcommand{\bkLgbmLiqT}{0.82}  % lgbm + liquidity vs HAR + cal.: HAC t
\newcommand{\bkLgbmMewSh}{1.16}  % lgbm + market_ew: Sharpe mid
\newcommand{\bkLgbmMewT}{0.97}  % lgbm + market_ew vs HAR + cal.: HAC t
\newcommand{\bkTreeTradeSigN}{four}  % tree x bucket trade gains with HAC t > 1.96
\newcommand{\bkXgbMewStepmN}{two}  % block lengths at which StepM (5 %) names XGBoost + market_ew
\newcommand{\bkXgbMewStepmBlocks}{10 and 21}
\newcommand{\bkXgbMewStepmMiss}{5}  % block lengths at which it is not named
\newcommand{\bkLgbmMomDM}{\ensuremath{-}2.01}  % LightGBM + moments vs HAR + cal., DM QLIKE
\newcommand{\bkPoolAllDM}{\ensuremath{-}0.65}  % six-estimator pool, all features vs HAR + cal., DM QLIKE
\newcommand{\bkPoolAllSh}{1.67}
\newcommand{\bkPoolAllBaseSh}{1.22}
\newcommand{\bkPoolAllT}{1.00}
\newcommand{\bkPoolFiveLiqOutSh}{2.03}  % five-estimator pool, all features - liquidity: Sharpe mid
\newcommand{\bkPoolFiveAllSh}{1.64}  % ... all features: Sharpe mid
\newcommand{\bkPoolFiveLiqOutT}{2.73}  % ... HAC t
\newcommand{\bkPoolFiveLiqOutSpaMax}{0.028}  % SPA consistent p (test B, trade), largest over blocks
\newcommand{\bkPoolTreeMomT}{2.26}  % tree pool + moments vs HAR + cal.: HAC t
\newcommand{\bkPoolTreeMewT}{2.43}  % tree pool + market_ew vs HAR + cal.: HAC t
\newcommand{\bkPoolTreeAllT}{0.75}  % tree pool, all features: HAC t
\newcommand{\bkMcsN}{576}  % candidates of the grand MCS
\newcommand{\bkMcsSets}{30}  % grand MCS sets on QLIKE, winsorized QLIKE and trade (blocks x stats x levels)
\fi
\subsection{Bucket by Bucket: Do the Inputs Help Only Together?}
\label{sec:app_buckets}

\paragraph{Question.} Does each bucket of exogenous inputs add little to HAR + calendar on
its own, while all features together add a discernible amount? If so, a joint test against
\texttt{baseline} should name \texttt{all\_features} and no single bucket.

\paragraph{Design.} The \bkNBuckets{} buckets partition the \bkNSeries{} series of
\texttt{all\_features}: the E-mini's return moments (\bkSizeMom{} series), liquidity
(\bkSizeLiq, the two order-flow imbalances included), the equal- and value-weighted stock
cross-sections \texttt{market\_ew} and \texttt{market\_vw} (\bkSizeMew{} and \bkSizeMvw),
sentiment (\bkSizeSent), implied volatility (\bkSizeIv), volatility demand (\bkSizeVd) and the
FOMC calendar (\bkSizeFomc). Each bucket gives two designs, HAR + calendar plus the bucket
and all features minus the bucket; with the two anchors each estimator fits \bkNDesigns{}
designs for the 16:00 bar. The estimators are the per-bar ridge, lasso and elastic net,
pre-tuned (Section~\ref{sec:close_option_methods}), and LightGBM, XGBoost and the random
forest with their shipped settings, refit every session. The \bkNPools{} pools average each
design's daily loss and daily midpoint P\&L over the linear models, over the trees, over all
\bkNEst{} estimators, and over all but LightGBM. Every forecast is its plain back-transform
$f^{2}B$, on the same \bkDays{} days.

\paragraph{Tests.} Pairwise, a design is set against its reference by the Diebold--Mariano
statistic of QLIKE, of a winsorized QLIKE (the daily loss differentials clipped at their
extreme percentiles across the designs) and of the fit-space squared error, and on the trade
by the $\mathrm{sign}(s)$ Sharpe ratio at the midpoint and the Newey--West $t$ of the daily
P\&L difference (significance at \bkSigPct\,\%). Jointly, Hansen's SPA test and
Romano--Wolf's StepM (circular block bootstrap, \bkBoot{} draws, blocks of \bkBlockA,
\bkBlockB{} and \bkBlockC{} days) ask on QLIKE, winsorized QLIKE and the trade loss:
(A) does a bucket alone, or all features, beat HAR + calendar; (B) does removing a bucket
from all features help; (C) does removing it hurt? A design is named when StepM rejects at a
family-wise rate of \bkFwer\,\% at every block length.

\paragraph{Finding.} The story fails in each of the \bkCells{} family-and-loss cells
(Table~\ref{tab:app_buckets}). StepM never names all features against HAR + calendar, at any
block length, and pairwise its gain is never significant: its most negative DM on QLIKE is
\bkAllDMMin{} (\bkAllDMMinWho) and its largest HAC $t$ on the trade \bkAllTMax{}
(\bkAllTMaxWho), although it trades above HAR + calendar in \bkAllShHigher{} of the
\bkNFam{} families. Where a gain is significant, it comes from one bucket. On QLIKE, for the
linear models, it is the value-weighted cross-section: StepM names \texttt{market\_vw} in
\bkBucketCells{} cells (\bkBucketCellsWho); with it the lasso's QLIKE changes by
\bkLassoVwDQ\,\% (DM \bkLassoVwDM, SPA $p\le\bkLassoVwSpaMax$), and the linear pool's DM is
\bkPoolLinVwDM{} (\bkPoolLinVwDMw{} winsorized) against \bkPoolLinAllDM{}
(\bkPoolLinAllDMw) for all features. Removing it is the removal that hurts: test C names it
in \bkRemoveCells{} cells (linear pool DM \bkPoolLinVwOutDM, \bkPoolLinVwOutDMw{}
winsorized). Liquidity added alone makes the linear models worse: QLIKE rises by
\bkLiqLinDQMin--\bkLiqLinDQMax\,\% (DM \bkLiqLinDMMin{} to \bkLiqLinDMMax). On the trade, for
the trees, the gains are single buckets and pairwise only. XGBoost trades at \bkXgbBaseSh{}
on HAR + calendar, at \bkXgbAllSh{} on all features (HAC $t$ \bkXgbAllT), at \bkXgbLiqSh{}
with liquidity alone ($t$ \bkXgbLiqT) and at \bkXgbMewSh{} with \texttt{market\_ew} alone
($t$ \bkXgbMewT); StepM names the latter at blocks of \bkXgbMewStepmBlocks{} days, not at
\bkXgbMewStepmMiss. The random forest moves from \bkRfBaseSh{} to \bkRfAllSh{} ($t$
\bkRfAllT) with all features, and to \bkRfMomSh{} ($t$ \bkRfMomT) and \bkRfMewSh{} ($t$
\bkRfMewT) with the moments or \texttt{market\_ew} alone. LightGBM has no StepM rejection;
its one significant pairwise gain is the moments on QLIKE (DM \bkLgbmMomDM). Pooled over all
\bkNEst{} estimators, all features against HAR + calendar has DM \bkPoolAllDM{} on QLIKE and
trades at \bkPoolAllSh{} against \bkPoolAllBaseSh{} ($t$ \bkPoolAllT). In the pool without
LightGBM, removing liquidity from all features raises the Sharpe ratio from
\bkPoolFiveAllSh{} to \bkPoolFiveLiqOutSh{} ($t$ \bkPoolFiveLiqOutT; test B at every block
length, SPA $p\le\bkPoolFiveLiqOutSpaMax$).

\paragraph{Reading.} On these \bkDays{} days the inputs do not help only together. No joint
test separates all features from HAR + calendar, and the single buckets that do separate
differ with the model and the loss: the value-weighted cross-section for the linear models'
accuracy; the moments, the equal-weighted cross-section and liquidity for the trees' trade,
where all features trades below each tree's best single bucket. The model confidence set
over the same study's \bkMcsN{} single-series candidates keeps every HAR + calendar model in
all \bkMcsSets{} of its sets on these losses, so it cannot rank the input sets either. The
families share designs and days, so their counts are not independent evidence.

\begin{table}[H]
\centering
\caption{The bucket story check by family, on \bkDays{} days. QLIKE: all features against
HAR + calendar (change of the mean in \%, DM; negative = all features better) and the
add-bucket design with the most negative DM. Trade: the $\mathrm{sign}(s)$ Sharpe ratio at
the midpoint of HAR + calendar, of all features (in parentheses the HAC $t$ of the
difference to HAR + calendar) and of the add-bucket design with the largest HAC $t$.
StepM: designs named at a family-wise rate of \bkFwer\,\% at every block length --- A, the
bucket added to HAR + calendar beats it; B, removing it from all features helps; C, removing
it hurts --- on QLIKE (Q), winsorized QLIKE (Q$_w$) or the trade loss (T); all features is
never named. Pools average each design's daily loss and P\&L over their estimators.}
\label{tab:app_buckets}
\resizebox{\textwidth}{!}{\scriptsize\setlength{\tabcolsep}{2.4pt}% AUTO-GENERATED by writeup/make_buckets_tex.py -- do not edit.
% Source: results/feature_mcs_2026-10-08/buckets_pairwise.csv, buckets_stepm.csv, buckets_verdicts.csv
\begin{tabular}{lrrlrlll}
\toprule
 & \multicolumn{3}{c}{QLIKE} & \multicolumn{3}{c}{trade: Sharpe (HAC $t$)} & \\
\cmidrule(lr){2-4}\cmidrule(lr){5-7}
family & \shortstack{all features\\$\Delta$\,\%} & \shortstack{all features\\DM} & \shortstack{best bucket\\(DM)} & \shortstack{HAR\\+ cal.} & \shortstack{all\\features} & \shortstack{best bucket} & \shortstack[l]{StepM, every\\block length} \\
\midrule
elastic net & \ensuremath{-}0.4 & \ensuremath{-}0.10 & \texttt{market\_vw} \ensuremath{-}1.80 & 1.37 & 1.92 (1.34) & \texttt{market\_vw} 1.90 (1.69) & C: \texttt{market\_vw} (Q$_w$) \\
ridge & 1.2 & 0.26 & \texttt{market\_vw} \ensuremath{-}1.65 & 1.37 & 1.42 (0.09) & \texttt{moments} 1.70 (0.77) & A: \texttt{market\_vw} (Q$_w$) \\
lasso & \ensuremath{-}3.4 & \ensuremath{-}1.22 & \texttt{market\_vw} \ensuremath{-}3.75 & 1.37 & 1.88 (1.47) & \texttt{market\_vw} 1.74 (1.58) & \begin{tabular}[t]{@{}l@{}}A: \texttt{market\_vw} (Q, Q$_w$)\\C: \texttt{market\_vw} (Q$_w$)\end{tabular} \\
LightGBM & \ensuremath{-}4.7 & \ensuremath{-}1.43 & \texttt{moments} \ensuremath{-}2.01 & 0.79 & 1.33 (1.11) & \texttt{moments} 1.61 (1.48) & --- \\
XGBoost & \ensuremath{-}4.2 & \ensuremath{-}1.41 & \texttt{moments} \ensuremath{-}1.68 & 0.95 & 1.12 (0.30) & \texttt{market\_ew} 1.65 (2.39) & --- \\
random forest & \ensuremath{-}0.2 & \ensuremath{-}0.05 & \texttt{moments} \ensuremath{-}1.11 & 0.66 & 0.85 (0.34) & \texttt{market\_ew} 1.47 (2.38) & --- \\
\midrule
linear pool & \ensuremath{-}0.9 & \ensuremath{-}0.23 & \texttt{market\_vw} \ensuremath{-}2.35 & 1.37 & 1.86 (0.97) & \texttt{market\_vw} 1.75 (1.38) & \begin{tabular}[t]{@{}l@{}}A: \texttt{market\_vw} (Q$_w$)\\C: \texttt{market\_vw} (Q$_w$)\end{tabular} \\
tree pool & \ensuremath{-}3.0 & \ensuremath{-}1.00 & \texttt{moments} \ensuremath{-}1.71 & 0.87 & 1.21 (0.75) & \texttt{market\_ew} 1.61 (2.43) & --- \\
pool of all six & \ensuremath{-}2.0 & \ensuremath{-}0.65 & \texttt{market\_vw} \ensuremath{-}1.80 & 1.22 & 1.67 (1.00) & \texttt{market\_vw} 1.69 (1.95) & --- \\
pool, no LightGBM & \ensuremath{-}1.4 & \ensuremath{-}0.45 & \texttt{market\_vw} \ensuremath{-}2.00 & 1.26 & 1.64 (0.88) & \texttt{market\_vw} 1.76 (1.88) & B: \texttt{liquidity} (T) \\
\bottomrule
\end{tabular}
}
\end{table}
